% ============================================================
%  Code listing style — fvextra line numbers + codeOS window frame
%
%  Included from preamble.tex via \input.
%  Depends on: xcolor (accent colors), tcolorbox [skins,breakable]
%  — both loaded earlier in preamble.tex.
% ============================================================

% ── Line numbers + soft wrap (fvextra) ───────────────────────────
% fvextra patches fancyvrb to support breaklines and line numbers.
% Must be loaded after fancyvrb (which pandoc's template loads first).
\usepackage{fvextra}

% Line number style: small, muted gray
\renewcommand{\theFancyVerbLine}{%
  \small\textcolor{textmuted}{\arabic{FancyVerbLine}}%
}

% Global defaults for all fancyvrb Verbatim environments
% (including pandoc's Highlighting environment).
\fvset{
  breaklines    = true,
  breakanywhere = false,
  numbers       = left,
  numbersep     = 12pt,
  xleftmargin   = 2em,
}

% ── codeOS window frame ───────────────────────────────────────────
% pandoc wraps syntax-highlighted blocks in a Shaded environment.
% We redefine it as a tcolorbox with a title-bar containing three
% traffic-light dots (left) and an optional centered filename.

% Colors
\definecolor{shadecolor}{HTML}{ffffff}    % fallback for framed/snugshade
\definecolor{codeosred}{HTML}{FF5F57}      % traffic-light red
\definecolor{codeosyellow}{HTML}{FEBC2E}   % traffic-light yellow
\definecolor{codeosgreen}{HTML}{28C840}    % traffic-light green
\definecolor{codeoswinframe}{HTML}{DEDEDE} % border + title separator
\definecolor{codewinbg}{HTML}{FFFFFF}    % window background

% Filename command — injected by filters/preamble.lua before each code
% block that carries a filename= attribute; cleared after each block.
\newcommand{\currentcodefilename}{}
\newcommand{\codeblockfilename}[1]{\gdef\currentcodefilename{#1}}

% Invisible right-side balancer so the filename stays centred in the bar.
\newcommand{\codeosphantom}{%
  \phantom{\tikz[baseline=-0.6ex, inner sep=0pt]{%
    \fill (0,0) circle[radius=5pt];%
    \fill (16pt,0) circle[radius=5pt];%
    \fill (32pt,0) circle[radius=5pt];%
  }}%
}

\renewenvironment{Shaded}{%
  \begin{tcolorbox}[%
    enhanced,
    breakable,
    colback      = codewinbg,
    colframe     = codeoswinframe,
    colbacktitle = codewinbg,
    coltitle     = codewinbg,
    boxrule      = 0.8pt,
    titlerule    = 0.5pt,
    arc          = 10pt,
    toptitle     = 8pt,
    bottomtitle  = 8pt,
    top          = 6pt,
    bottom       = 6pt,
    left         = 6pt,
    right        = 6pt,
    title        = {%
      % Traffic-light dots — zero-width anchor so \hfill centres the filename
      \makebox[0pt][l]{%
        \tikz[baseline=-0.6ex, inner sep=0pt]{%
          \fill[codeosred]    (0,0)    circle[radius=5pt];%
          \fill[codeosyellow] (16pt,0) circle[radius=5pt];%
          \fill[codeosgreen]  (32pt,0) circle[radius=5pt];%
        }%
      }%
      % Filename centred via \hfill + phantom mirror on the right
      \hfill
      \ifx\currentcodefilename\empty\else
        {\ttfamily\footnotesize\color{textmuted}\currentcodefilename}%
      \fi
      \hfill\codeosphantom
    },
  ]%
}{%
  \end{tcolorbox}%
  \gdef\currentcodefilename{}%
}
